\documentclass[11pt,a4paper]{article}

% ============================================================
% Packages (core only — texlive-latex-base + texlive-base)
% ============================================================
\usepackage{amsmath, amssymb, amsthm}
\usepackage[a4paper, margin=2.5cm]{geometry}
\usepackage{array}

% ============================================================
% Theorem environments
% ============================================================
\theoremstyle{definition}
\newtheorem{definition}{Definition}[section]
\newtheorem{remark}{Remark}[section]
\newtheorem{proposition}{Proposition}[section]

\theoremstyle{plain}
\newtheorem{theorem}{Theorem}[section]
\newtheorem{lemma}{Lemma}[section]

% ============================================================
% Notation macros
% ============================================================
\newcommand{\R}{\mathbb{R}}
\newcommand{\E}{\mathbb{E}}
\newcommand{\N}{\mathbb{N}}
\newcommand{\ind}{\mathbf{1}}
\newcommand{\norm}[1]{\left\lVert #1 \right\rVert}
\newcommand{\inner}[2]{\langle #1,\, #2 \rangle}
\newcommand{\corr}{\operatorname{corr}}
\newcommand{\sgn}{\operatorname{sign}}
\newcommand{\code}[1]{\texttt{\detokenize{#1}}}
\DeclareMathOperator*{\argmax}{arg\,max}

% ============================================================
\begin{document}

\title{\textbf{From Liquidation Geometry to Price Movement}\\[4pt]
\large A Spatial-Density Heatmap and Cascade-Probability Model\\
for Cryptocurrency Perpetual Futures}

\author{Yigit Onat\\
\textit{Independent Researcher}\\
\texttt{yionat@gmail.com}}
\date{June 2026}
\maketitle

% ============================================================
\begin{abstract}
Leveraged perpetual-futures positions carry deterministic liquidation prices; when price reaches a
dense cluster of them, forced market liquidations impose price impact that can reach the next
cluster, producing a self-reinforcing cascade. We formalize this intuition. We construct a
\emph{spatial liquidation-density field}---a ``liquidity heatmap''---that discretizes the
log-return axis around the mid-price into $K=200$ bins ($10$ basis points each, $\pm 10\%$) and
maps the USD notional of each tracked position to the bin containing its liquidation price. From
this field we extract eight scale- and level-invariant geometric features (nearest-cluster
distance, mass concentration, cascade-chain length, directional asymmetry, absorption capacity,
cluster velocity, mass-weighted distance, and spatial entropy). We then give a structural
price-movement model in which a closed-form domino condition,
$\delta_B-\delta_A < \lambda\,M_A/m_t$, determines whether liquidating one cluster triggers the
next, with order-book depth acting as a damper (a liquidity spiral) and the directional asymmetry
of the field predicting cascade sign. The geometry feeds an online, delayed-target logistic model
of cascade probability. The contribution of this paper is the construction and the model, together
with an explicit, falsifiable five-gate validation protocol and a candid treatment of the binding
practical obstacle: on a decentralized venue the field is observed only through a \emph{sample} of
tracked wallets, not the full population of open interest. We formalize this sample-to-population
bias, propose an open-interest rescaling that preserves field \emph{shape}, and report the failure
mode it produces in practice. We deliberately present no fabricated performance numbers; empirical
validation under the stated protocol is the explicit next step.
\end{abstract}

\smallskip
\noindent\textbf{Keywords:} market microstructure, liquidation cascade, price impact, perpetual
futures, spatial density estimation, fire sales, order-book liquidity, cryptocurrency

\smallskip
\noindent\textbf{JEL Classification:} G12, G14, G17, C53

\tableofcontents
\bigskip

% ============================================================
\section{Introduction}
\label{sec:intro}
% ============================================================

\subsection{The problem}

Perpetual futures are the dominant instrument in cryptocurrency derivatives, and they are traded
with high leverage. A leveraged position has a deterministic \emph{liquidation price}: the level at
which the maintenance-margin constraint binds and the exchange force-closes the position with a
market order. These forced orders are price-takers and price-movers. When the mid-price approaches
a region of the price axis that is dense with liquidation prices, a move into that region triggers
a burst of same-signed market orders whose impact can carry the price into the \emph{next} dense
region---a liquidation cascade, the cryptocurrency analogue of a fire sale.

Two facts make this mechanism tractable as a microstructure object. First, the trigger levels are
not latent: a position's liquidation price is a deterministic function of its entry, leverage, and
side. Second, the impact channel is the classical one: forced liquidations consume order-book
depth, and their price impact is governed by market illiquidity. The open question is therefore not
\emph{whether} liquidation geometry matters, but how to \emph{represent} it as a field over the
price axis and how that field's geometry maps to future price movement.

\subsection{Approach and contributions}

We represent the liquidation landscape as a one-dimensional density field $H(t,\cdot)$ over the
log-return axis relative to the current mid-price, updated at $100\,$ms resolution
(Section~\ref{sec:heatmap}). From the field we extract eight geometric features that are invariant
to price level and, being ratios or distances, approximately invariant to the unobserved scale of
total open interest (Section~\ref{sec:features}). We then give a structural model of how the field
moves price (Section~\ref{sec:cascade}): a single cluster's impact follows Kyle (1985), a
closed-form condition determines whether it triggers the next cluster, order-book depth damps the
chain in the manner of Brunnermeier and Pedersen (2009), and the field's directional asymmetry
predicts the sign. The geometry drives an online logistic model of cascade probability with delayed
targets (Section~\ref{sec:probmodel}), gated by an explicit five-criterion validation protocol
(Section~\ref{sec:validation}).

Specifically, this paper contributes:
\begin{enumerate}
  \item A construction of the liquidation-density field in \emph{log-return} coordinates, making
        bins symmetric in percentage terms and invariant to price level, with an $O(N_t)$ update.
  \item Eight scale-invariant geometric features of the field, each with a microstructural
        interpretation.
  \item A structural price-movement model with a closed-form domino-trigger condition and an
        order-book absorption damper, linking field geometry to cascade magnitude and direction.
  \item An online, delayed-target cascade-probability model and a falsifiable five-gate validation
        protocol.
  \item A formal treatment of the sample-to-population observation problem on decentralized venues,
        an open-interest rescaling that preserves field shape, and an honest account of the
        coverage-driven failure mode observed in practice.
\end{enumerate}

We emphasize at the outset that this is a \emph{methods and model} paper. The geometry and the model
are fully specified and implemented; a full out-of-sample evaluation under the protocol of
Section~\ref{sec:validation} is stated as future work, and we report no measured performance figures
we cannot substantiate.

% ============================================================
\section{Related Work}
\label{sec:related}
% ============================================================

\paragraph{Price impact of forced trades.}
Kyle (1985) gives the canonical linear price-impact model in which order flow $q$ moves price by
$\lambda q$; $\lambda$ is the inverse market depth. Forced liquidations are exactly such order flow.
Our single-cluster impact and the domino condition are built directly on this primitive, with
$\lambda$ estimated from recent trades in the spirit of Amihud (2002) and Hasbrouck (2009).

\paragraph{Fire sales and endogenous risk.}
The cascade mechanism is a fire sale in the sense of Caccioli, Shrestha, Moore, and Farmer (2014),
who analyze contagion in overlapping portfolios, and of Cont and Wagalath (2016), who measure
endogenous risk created by feedback between forced selling and price. Our cascade-chain feature is
the length of the triggered component of exactly such a contagion graph.

\paragraph{Liquidity spirals.}
Brunnermeier and Pedersen (2009) show how market and funding liquidity interact to amplify shocks
in thin markets. Our absorption term---order-book depth divided by cluster mass---is the
discrete, level-local form of their damper: a deep book absorbs a liquidation without propagation,
a thin one lets price gap through.

\paragraph{Cryptocurrency liquidations.}
Perpetual-futures liquidation cascades are a recognized empirical phenomenon in crypto markets, but
they are typically studied ex post from exchange liquidation feeds. Our contribution is a
\emph{forward-looking, spatial} representation built from position data, and an explicit treatment
of the fact that on a decentralized venue this data is sampled, not complete.

% ============================================================
\section{The Liquidity Heatmap: Construction}
\label{sec:heatmap}
% ============================================================

\subsection{Inputs and coordinates}

Let $m_t$ denote the mid-price at time $t$. The system tracks a set of leveraged positions; position
$i$ has notional $V_i$ (USD), liquidation price $\ell_i$, side, and entry price. We place each
position on the \emph{log-return axis} relative to the mid,
\begin{equation}
  \delta_i = \ln\!\frac{\ell_i}{m_t},
  \label{eq:logreturn}
\end{equation}
so that $\delta_i$ is the (signed) percentage distance of the liquidation level from the current
price. The axis is discretized into $K$ uniform bins on $[-R, R]$:
\begin{equation}
  \delta_k = -R + \big(k+\tfrac12\big)\Delta,\quad
  \Delta = \frac{2R}{K},\quad
  \mathcal{B}_k = \big(-R + k\Delta,\; -R + (k{+}1)\Delta\big),\quad k = 0,\dots,K-1,
  \label{eq:bins}
\end{equation}
with defaults $K = 200$, $R = 0.10$ (so $\Delta = 0.001$, i.e.\ $10$ basis points per bin,
spanning $\pm 10\%$). Working in log-returns makes the bins symmetric in percentage terms and
independent of the absolute price level.

\subsection{The density field}

\begin{definition}[Liquidation-density field]
\label{def:field}
The liquidity heatmap is the field
\begin{equation}
  H(t,k) = \sum_{i=1}^{N_t} V_i \,\ind\!\big[\delta_i \in \mathcal{B}_k\big]\,
           \ind\!\big[V_i \geq \tau_V\big],
  \qquad k = \Big\lfloor \tfrac{\delta_i + R}{\Delta} \Big\rfloor,
  \label{eq:field}
\end{equation}
i.e.\ $H(t,k)$ is the total USD notional whose liquidation price lies in bin $k$. Positions with
notional below $\tau_V$ (default $\$1{,}000$) or with $\delta_i \notin [-R,R]$ are discarded.
\end{definition}

The field is rebuilt in full each snapshot; each position maps to one bin via a logarithm, an
addition, and a floor, so the update is $O(N_t)$. There is no temporal smoothing of $H$ itself;
temporal information enters only through the velocity feature below.

\subsection{Derived channels}

Three further channels give the field local structure (here $k_0 = K/2$ is the at-the-money bin):
\begin{align}
  \text{cumulative mass:} \quad & C(t,k) = \textstyle\sum_{j \text{ between } k \text{ and } k_0} H(t,j),
    \label{eq:cumulative}\\
  \text{cluster gradient:} \quad & G(t,k) = H(t,k) - H(t,k-1),
    \label{eq:gradient}\\
  \text{temporal change:} \quad & T(t,k) = \alpha_1 \frac{H(t,k)-H(t-\Delta t_1,k)}{\Delta t_1}
                                          + \alpha_2 \frac{H(t,k)-H(t-\Delta t_2,k)}{\Delta t_2},
    \label{eq:temporal}
\end{align}
with $\Delta t_1 = 600$ ticks ($1$\,min), $\Delta t_2 = 3000$ ticks ($5$\,min), and weights
$\alpha_1 = 0.7$, $\alpha_2 = 0.3$. The cumulative channel measures the fuel between the mid and a
level (cascade reach), the gradient locates cluster boundaries, and the temporal channel is the
cluster-migration velocity field.

% ============================================================
\section{Spatial Features}
\label{sec:features}
% ============================================================

Eight scalar features summarize the field's geometry each snapshot. Write the total mass
$\Sigma_H = \sum_k H(t,k)$, $\bar H = \Sigma_H/K$, and let $p_k = H(t,k)/\Sigma_H$ be the
normalized mass. Let $\tau_{\text{cl}} = \$10^6$ and $\tau_{\text{ch}} = \$10^5$ be the cluster and
chain thresholds.

\begin{table}[ht]
\begin{center}\small
\renewcommand{\arraystretch}{1.5}
\begin{tabular}{p{3.4cm} p{6.2cm} p{4.6cm}}
\hline
\textbf{Feature} & \textbf{Definition} & \textbf{Interpretation} \\
\hline
\code{hm_nearest_cluster_dist} &
  $d_{\min} = \min\{|\delta_k| : H(t,k) > \tau_{\text{cl}}\}$ &
  distance to the nearest large cluster (smaller $=$ cascade closer) \\
\code{hm_cluster_mass_ratio} &
  $\mathrm{CMR} = H(t,k^\ast)/\bar H$, $k^\ast = \argmax_k H$ &
  peak concentration relative to mean \\
\code{hm_cascade_chain_length} &
  longest run of consecutive bins with $H > \tau_{\text{ch}}$ within $|\delta|\le 0.05$ &
  domino potential (triggered-chain length) \\
\code{hm_asymmetric_cascade_pot} &
  $\mathrm{ACP} = \dfrac{\sum_{0<\delta\le.05}H - \sum_{-.05\le\delta<0}H}{\sum_{|\delta|\le.05}H + \epsilon}$ &
  directional bias: $>0$ short-squeeze, $<0$ long-cascade \\
\code{hm_absorption_capacity} &
  $A = \min\!\big(\tfrac{D^{\text{bid}}(.02)}{\text{down mass}+\epsilon},\,
       \tfrac{D^{\text{ask}}(.02)}{\text{up mass}+\epsilon}\big)$ &
  book depth per unit liquidation mass ($<1$ $=$ insufficient damping) \\
\code{hm_cluster_velocity} &
  $v = \big(d_{\min}(t) - d_{\min}(t-\tau)\big)/\tau$, $\tau = 600$ &
  rate of approach of the nearest cluster ($<0$ approaching) \\
\code{hm_mass_weighted_distance} &
  $\bar d = \dfrac{\sum_k |\delta_k|\,H(t,k)}{\Sigma_H + \epsilon}$ &
  centre of liquidation gravity \\
\code{hm_heatmap_entropy} &
  $S = -\sum_k p_k \ln p_k$ &
  spatial concentration ($S$ low $=$ tight clusters) \\
\hline
\end{tabular}
\end{center}
\caption{The eight liquidity-heatmap features. $\epsilon$ is a small constant ($1$ USD).
Absorption uses order-book depth within $2\%$ on each side; when a side carries no mass its ratio is
treated as $+\infty$, and $A$ is clamped to $[0,100]$.}
\label{tab:features}
\end{table}

These features are deliberately constructed to be \emph{ratios or signed distances}
(Table~\ref{tab:features}): $\mathrm{CMR}$, $\mathrm{ACP}$, $S$, $\bar d$, and $d_{\min}$ are
unchanged if the entire field is rescaled by a common factor. This is the property that makes them
usable when only a fraction of total open interest is observed (Section~\ref{sec:data}).

% ============================================================
\section{The Cascade Mechanism: A Price-Movement Model}
\label{sec:cascade}
% ============================================================

\subsection{Single-cluster impact}

Let cluster $A$ at $\delta_A$ carry mass $M_A = H(t,k_A)$. Treating its forced liquidation as an
order of notional $M_A$ and applying Kyle (1985) with depth parameter $\lambda$ (estimated from
recent trades), its price impact in log-return units is
\begin{equation}
  \Delta p_A = \lambda \,\frac{M_A}{m_t}.
  \label{eq:impact}
\end{equation}

\subsection{The domino condition}

A second cluster $B$ at $\delta_B > \delta_A$ (on the same side of the mid) is reached if and only
if cluster $A$'s impact closes the gap:
\begin{equation}
  \boxed{\;\delta_B - \delta_A < \lambda\,\frac{M_A}{m_t}\;}
  \label{eq:domino}
\end{equation}
Equation~\eqref{eq:domino} is the structural heart of the model. It depends on the \emph{spacing}
of clusters (captured by the gradient channel and the chain-length feature) and their \emph{mass}
(the concentration feature), mediated by illiquidity $\lambda$. When it holds across a run of
clusters $\delta_1 < \dots < \delta_n$, the total move is the sum of triggered masses,
\begin{equation}
  \Delta p_{\text{total}} \approx \frac{\lambda}{m_t} \sum_{j:\,\text{triggered}} M_j,
  \label{eq:total}
\end{equation}
and the number of triggered nodes is exactly the cascade-chain feature, read off a Caccioli
et al.\ (2014) contagion graph.

\subsection{Absorption damping}

Order-book depth opposes the cascade. With bid depth $D^{\text{bid}}(\delta_A)$ available to absorb
a downward liquidation at $A$, define the damping
\begin{equation}
  \eta = \min\!\Big(1,\; \frac{D^{\text{bid}}(\delta_A)}{M_A}\Big),
  \qquad
  \Delta p_A^{\text{eff}} = (1-\eta)\,\Delta p_A.
  \label{eq:damping}
\end{equation}
When $\eta \to 1$ (a deep ``wall,'' i.e.\ absorption $A \gg 1$) the impact is absorbed and the chain
breaks; when $\eta \to 0$ (a thin book) it is amplified---the Brunnermeier--Pedersen (2009) spiral.
This is the mechanism behind support/resistance walls and absorption.

\begin{proposition}[Magnet and rejection]
\label{prop:magnet}
Under \eqref{eq:domino}--\eqref{eq:damping}, a large cluster within reach acts as a price
\emph{magnet}: conditional on the mid entering its neighborhood, the probability of a move toward
it increases with $M_A$ and decreases with the local absorption $A$. Conversely, a cluster fronted
by a deep book ($\eta\approx 1$) produces \emph{rejection} rather than breakout.
\end{proposition}

\subsection{Direction}

The sign of the field's near-money asymmetry predicts the cascade direction. With $\mathrm{ACP}$ as
in Table~\ref{tab:features}, mass concentrated above the mid is short-liquidation fuel (an upward
squeeze) and mass below is long-liquidation fuel (a downward cascade), so the predicted direction
is $\sgn(\mathrm{ACP})$.

% ============================================================
\section{A Probabilistic Cascade Model}
\label{sec:probmodel}
% ============================================================

\subsection{Event label}

We define a cascade over horizon $T$ as a large \emph{liquidation-driven} move---a price move
conjoined with abnormal liquidation volume, which distinguishes it from an ordinary news move:
\begin{equation}
  y_t = \ind\!\Big[\,|\Delta p_{t\to t+T}| > X\% \;\wedge\; \mathrm{LiqVol}_{t\to t+T} > Y\,\Big],
  \qquad \Delta p = \ln\frac{m_{t+T}}{m_t},
  \label{eq:label}
\end{equation}
with defaults $X = 3\%$, $Y = \$500\text{K}$, $T = 3000$ ticks ($5$\,min), giving a base rate near
$5\%$. Because per-tick liquidation volume is not directly observable on the venue, the
implementation currently uses the price-threshold component as a proxy---a conservative choice we
flag explicitly.

\subsection{Features and model}

The model consumes an $11$-dimensional vector: the eight features of Table~\ref{tab:features} plus
three interactions chosen from the mechanism,
\begin{equation}
  \mathbf{x}(t) = \big[\,d_{\min},\,\mathrm{CMR},\,L,\,\mathrm{ACP},\,A^{-1},\,v,\,\bar d,\,S,\;
                       \underbrace{L\!\cdot\!A^{-1}}_{\text{fuel, no damper}},\;
                       \underbrace{v\!\cdot\!\mathcal{I}}_{\text{approach}\times\text{illiquidity}},\;
                       \underbrace{\mathrm{ACP}\!\cdot\!\sgn(f)}_{\text{aligned with funding}}\,\big]^\top,
  \label{eq:xvec}
\end{equation}
where $L$ is the chain length, $\mathcal{I}$ a composite illiquidity, and $f$ the funding rate.
Cascade probability is an online logistic regression,
\begin{equation}
  P(\text{cascade}\mid \mathbf{x}) = \sigma\!\big(\beta_0 + \boldsymbol\beta^\top \mathbf{x}\big),
  \qquad \sigma(z) = \frac{1}{1+e^{-z}},
  \label{eq:logistic}
\end{equation}
trained by stochastic gradient descent on the cross-entropy loss with $\ell_2$ regularization
($\lambda_{\text{reg}} = 10^{-3}$). With per-step gradient $g_t = \hat p_t - y_t$,
\begin{equation}
  \boldsymbol\beta \leftarrow \boldsymbol\beta - \eta_t\big(g_t\,\mathbf{x}_t + \lambda_{\text{reg}}\boldsymbol\beta\big),
  \qquad \beta_0 \leftarrow \beta_0 - \eta_t g_t,
  \qquad \eta_t = \eta_0\,\gamma^t,
  \label{eq:sgd}
\end{equation}
with $\eta_0 = 0.01$, $\gamma = 0.9999$. Because the label \eqref{eq:label} is only known $T$ ticks
later, updates are applied from a \emph{delayed-target} buffer once each observation matures; inputs
are standardized online by Welford's algorithm, and the logit is clipped to $[-20,20]$. A logistic
model is deliberate: cascades are rare ($\approx 5\%$), leverage regimes are non-stationary, and the
coefficients $\boldsymbol\beta$ are interpretable as the weight the data place on each geometric
feature.

% ============================================================
\section{Validation Protocol}
\label{sec:validation}
% ============================================================

The model is only credible if it clears pre-registered gates. We specify five; they are
\emph{thresholds to be met}, not results claimed here.

\begin{center}\small
\begin{tabular}{l l}
\hline
\textbf{Gate} & \textbf{Criterion} \\
\hline
G1 — Discrimination & out-of-sample $\mathrm{AUC} > 0.65$ \emph{and} top-decile lift $> 2\times$ \\
G2 — Net information & $\mathrm{IC}_{\text{net}} = \mathrm{IC}_{\text{gross}} - \tfrac{2\bar s}{\sigma_r\sqrt T} > 0.02$ \\
G3 — Stability & walk-forward median $\mathrm{AUC} > 0.60$, $<20\%$ of windows below $0.55$ \\
G4 — Universality & BTC, ETH, SOL each pass G1, and $\cos(\boldsymbol\beta_i,\boldsymbol\beta_j) > 0.6$ pairwise \\
G5 — Orthogonality & $R^2 < 0.30$ regressing $\hat p$ on a spread/depth composite \\
\hline
\end{tabular}
\end{center}

\noindent Here $\bar s$ is the cascade-event spread (a cost proxy) and $\sigma_r$ the return
volatility. G2 ensures the signal survives transaction costs; G5 ensures it is not a relabeling of
ordinary liquidity. The prior empirical anchor that motivates the whole construction is the
hypothesis-test result that coarse liquidation-cluster proximity predicts large moves with lift
$> 2\times$; the heatmap's claim is that its eight \emph{spatial} features carry information
\emph{incremental} to that coarse proximity---a claim Gate~G5 and an explicit ablation against
\code{hm_nearest_cluster_dist} alone are designed to test.

% ============================================================
\section{Data Availability and the Sample-to-Population Problem}
\label{sec:data}
% ============================================================

The model's binding obstacle is observational, not theoretical. On a decentralized venue, position
data is obtained by querying the clearinghouse state of \emph{known wallet addresses}; no endpoint
returns the full population of open positions. Wallets are discovered from the public trade stream
and polled, yielding coverage on the order of $50$--$500$ wallets per symbol---an \emph{unknown}
fraction of open interest. The field of Definition~\ref{def:field} is therefore a sample,
$H^{\text{obs}}(t,k) = \rho_t\, H(t,k) + \text{noise}$, with sampling rate $\rho_t$ unobserved and
time-varying.

\paragraph{What survives sampling.} The scale-invariant features---$d_{\min}$, $\mathrm{CMR}$,
$\mathrm{ACP}$, $\bar d$, $S$---depend on the field's \emph{shape}, which is preserved in
expectation under uniform sampling; and the largest, cascade-causing positions are
disproportionately likely to be discovered (they dominate trade flow). What is biased are the
\emph{absolute-mass} quantities: absorption capacity $A$ is overestimated (observed mass
underestimates the true denominator), and chain length may be undercounted when untracked positions
fill gaps in \eqref{eq:domino}.

\paragraph{A shape-preserving rescaling.} Because total open interest \emph{is} observable, we can
restore scale without distorting shape,
\begin{equation}
  H^{\text{scaled}}(t,k) = H^{\text{obs}}(t,k)\cdot
       \frac{\mathrm{OI}_{\text{total}}(t)}{\sum_k H^{\text{obs}}(t,k)},
  \label{eq:oiscale}
\end{equation}
which leaves every ratio feature unchanged while correcting the mass-dependent ones. Validating
\eqref{eq:oiscale} against realized cascade magnitudes is itself a clean empirical sub-study.

\paragraph{Observed failure mode.} We report candidly that, in a preliminary multi-symbol backtest,
the cascade-probability algorithm produced \emph{zero} trades and the nearest-cluster-distance
feature frequently returned its no-data sentinel---direct evidence that, under current wallet
coverage, the observed field is too sparse for the absolute-mass features to fire. This is the
single most important practical result of the work to date: the geometry and the model are sound,
but their utility is gated on coverage, and the first priority is therefore data---aggressive wallet
discovery and the rescaling \eqref{eq:oiscale}---not further modeling.

% ============================================================
\section{Discussion}
\label{sec:discussion}
% ============================================================

The model makes several falsifiable claims that a venue with adequate coverage could test directly.
The domino condition \eqref{eq:domino} predicts that the probability of a multi-cluster cascade
jumps discontinuously as inter-cluster spacing falls below $\lambda M_A/m_t$; the absorption
relation \eqref{eq:damping} predicts that deep near-money book depth converts approach into
rejection rather than breakout (Proposition~\ref{prop:magnet}); and the directional claim predicts
that $\sgn(\mathrm{ACP})$ leads cascade sign. Each is a clean hypothesis with a pre-specified
statistic. The most important open question is incremental information: do the eight spatial
features add predictive content beyond simple cluster proximity? The construction is designed so
that this question is answerable by ablation, and Gate~G5 guards against the field being a disguised
liquidity proxy.

The honest status is that this is a fully specified instrument awaiting the data to validate it.
That ordering---specify the mechanism, expose the observational constraint, fix the data, then
measure---is, we argue, the correct one for a phenomenon whose theoretical structure is clear but
whose measurement on a decentralized venue is genuinely hard.

% ============================================================
\section{Conclusion}
\label{sec:conclusion}
% ============================================================

We have constructed a spatial liquidation-density field for cryptocurrency perpetual futures,
extracted eight scale-invariant geometric features from it, and given a structural price-movement
model in which a closed-form domino condition, an order-book damper, and a directional asymmetry map
the field's geometry to cascade magnitude, occurrence, and sign. The geometry drives an online,
delayed-target cascade-probability model, gated by an explicit five-criterion protocol. We have been
deliberate in not reporting performance numbers we cannot substantiate, and in foregrounding the
binding constraint: on a decentralized venue the field is sampled, not observed in full, and the
immediate priority is the coverage and rescaling needed to make the absolute-mass features fire.
With that data in hand, every claim above becomes a pre-registered, testable hypothesis.

% ============================================================
\section*{References}
\addcontentsline{toc}{section}{References}
% ============================================================

\begin{enumerate}
  \item Amihud, Y.\ (2002). Illiquidity and stock returns: cross-section and time-series effects.
        \textit{Journal of Financial Markets}, 5(1), 31--56.

  \item Brunnermeier, M.K.\ \& Pedersen, L.H.\ (2009). Market liquidity and funding liquidity.
        \textit{Review of Financial Studies}, 22(6), 2201--2238.

  \item Caccioli, F., Shrestha, M., Moore, C.\ \& Farmer, J.D.\ (2014). Stability analysis of
        financial contagion due to overlapping portfolios. \textit{Journal of Banking \& Finance},
        46, 233--245.

  \item Cont, R.\ \& Wagalath, L.\ (2016). Fire sales forensics: measuring endogenous risk.
        \textit{Mathematical Finance}, 26(4), 835--866.

  \item Hasbrouck, J.\ (2009). Trading costs and returns for U.S.\ equities: estimating effective
        costs from daily data. \textit{Journal of Finance}, 64(3), 1445--1477.

  \item Kyle, A.S.\ (1985). Continuous auctions and insider trading. \textit{Econometrica}, 53(6),
        1315--1335.

  \item Roll, R.\ (1984). A simple implicit measure of the effective bid-ask spread in an efficient
        market. \textit{Journal of Finance}, 39(4), 1127--1139.

  \item Bandt, C.\ \& Pompe, B.\ (2002). Permutation entropy: a natural complexity measure for time
        series. \textit{Physical Review Letters}, 88(17), 174102.

  \item Welford, B.P.\ (1962). Note on a method for calculating corrected sums of squares and
        products. \textit{Technometrics}, 4(3), 419--420.
\end{enumerate}

\end{document}
